% GENERATED FILE. Edit canonical lessons, then run npm run generate:books.
% canonical-source-hash: fnv1a64:9f83c973817cba30
% canonical-lessons: UR-C65-patthar, UR-C65-ret, UR-C65-khet, UR-C65-sitara, UR-C65-badal

\chapter{Stone, Sand, Field, Star, Cloud}
\label{ch:ur-stone-sand-field-star-cloud}
\hlchaptermodality{65}

\begin{chapteropening}
\textbf{By the end of this chapter:} \emph{I can say a stone, sand, a field, a star and a cloud.}

Complete the last lesson of chapter 65: I can say a stone, sand, a field, a star and a cloud.
\end{chapteropening}

\section[\texorpdfstring{\ur{پتھر} (patthar)}{پتھر (patthar)}]{\ur{پتھر} (patthar) — a stone}

\label{lesson:UR-C65-patthar}



\begin{quote}
Before the new one: say the Urdu for a river, then the Urdu for a lake.
\end{quote}

\subsection*{You'll want to know: \ur{پتھر}}
\textbf{\ur{پتھر}} — \emph{patthar} — \textquotedblleft{}a stone\textquotedblright{}.

\begin{cousinweb}
\textbf{\ur{پتھر}} (\emph{patthar}) is a tadbhava, a word inherited from Sanskrit and reshaped by sound change. It comes from \emph{prastara}, \textquotedblleft{}anything strewn; flat surface; stone\textquotedblright{} (CDIAL 8857).
\end{cousinweb}

\begin{grammarlens}[title={What you've built}]
One more word for this chapter.
\end{grammarlens}

\subsection*{Your turn}
\begin{itemize}
\raggedright
  \item \emph{Say it:} \emph{patthar}
  \item \emph{Say it:} \emph{patthar}, once more
  \item \emph{Say it:} say \emph{patthar}
  \item \emph{Recall:} say the Urdu for a river, then the Urdu for a lake, then say \emph{patthar} again
\end{itemize}

\begin{tcolorbox}[breakable,colback=teal!4,colframe=teal!35!black,title={Before you move on}]
What does \ur{پتھر} mean? (\textquotedblleft{}A stone\textquotedblright{}.) Say it once more.
\end{tcolorbox}
\section[\texorpdfstring{\ur{ریت} (ret)}{ریت (ret)}]{\ur{ریت} (ret) — sand}

\label{lesson:UR-C65-ret}



\begin{quote}
Before the new one: say the Urdu for a lake, then the Urdu for a stone.
\end{quote}

\subsection*{You'll want to know: \ur{ریت}}
\textbf{\ur{ریت}} — \emph{ret} — \textquotedblleft{}sand\textquotedblright{}.

\begin{cousinweb}
\textbf{\ur{ریت}} (\emph{ret}) is a tadbhava, a word inherited from Sanskrit and reshaped by sound change. It comes from \emph{retra}, \textquotedblleft{}perfumed powder\textquotedblright{} (CDIAL 10816). CDIAL doubts a suggested Munda source; the Sanskrit word is known only from the old word-lists.
\end{cousinweb}

\begin{grammarlens}[title={What you've built}]
One more word for this chapter.
\end{grammarlens}

\subsection*{Your turn}
\begin{itemize}
\raggedright
  \item \emph{Say it:} \emph{ret}
  \item \emph{Say it:} \emph{ret}, once more
  \item \emph{Say it:} say \emph{ret}
  \item \emph{Recall:} say the Urdu for a lake, then the Urdu for a stone, then say \emph{ret} again
\end{itemize}

\begin{tcolorbox}[breakable,colback=teal!4,colframe=teal!35!black,title={Before you move on}]
What does \ur{ریت} mean? (\textquotedblleft{}Sand\textquotedblright{}.) Say it once more.
\end{tcolorbox}
\section[\texorpdfstring{\ur{کھیت} (khet)}{کھیت (khet)}]{\ur{کھیت} (khet) — a field}

\label{lesson:UR-C65-khet}



\begin{quote}
Before the new one: say the Urdu for a stone, then the Urdu for sand.
\end{quote}

\subsection*{You'll want to know: \ur{کھیت}}
\textbf{\ur{کھیت}} — \emph{khet} — \textquotedblleft{}a field\textquotedblright{}.

\begin{cousinweb}
\textbf{\ur{کھیت}} (\emph{khet}) is a tadbhava, a word inherited from Sanskrit and reshaped by sound change. It comes from \emph{kṣetra}, \textquotedblleft{}land, field\textquotedblright{} (CDIAL 3735).
\end{cousinweb}

\begin{grammarlens}[title={What you've built}]
One more word for this chapter.
\end{grammarlens}

\subsection*{Your turn}
\begin{itemize}
\raggedright
  \item \emph{Say it:} \emph{khet}
  \item \emph{Say it:} \emph{khet}, once more
  \item \emph{Say it:} say \emph{khet}
  \item \emph{Recall:} say the Urdu for a stone, then the Urdu for sand, then say \emph{khet} again
\end{itemize}

\begin{tcolorbox}[breakable,colback=teal!4,colframe=teal!35!black,title={Before you move on}]
What does \ur{کھیت} mean? (\textquotedblleft{}A field\textquotedblright{}.) Say it once more.
\end{tcolorbox}
\section[\texorpdfstring{\ur{ستارہ} (sitārā)}{ستارہ (sitārā)}]{\ur{ستارہ} (sitārā) — a star}

\label{lesson:UR-C65-sitara}



\begin{quote}
Before the new one: say the Urdu for sand, then the Urdu for a field.
\end{quote}

\subsection*{You'll want to know: \ur{ستارہ}}
\textbf{\ur{ستارہ}} — \emph{sitārā} — \textquotedblleft{}a star\textquotedblright{}.

\begin{cousinweb}
\textbf{\ur{ستارہ}} (\emph{sitārā}) is a loan, usually traced to Persian \emph{sitāra}, \textquotedblleft{}star\textquotedblright{}. It is usually taken as a distant cousin of English \emph{star}.
\end{cousinweb}

\begin{grammarlens}[title={What you've built}]
One more word for this chapter.
\end{grammarlens}

\subsection*{Your turn}
\begin{itemize}
\raggedright
  \item \emph{Say it:} \emph{sitārā}
  \item \emph{Say it:} \emph{sitārā}, once more
  \item \emph{Say it:} say \emph{sitārā}
  \item \emph{Recall:} say the Urdu for sand, then the Urdu for a field, then say \emph{sitārā} again
\end{itemize}

\begin{tcolorbox}[breakable,colback=teal!4,colframe=teal!35!black,title={Before you move on}]
What does \ur{ستارہ} mean? (\textquotedblleft{}A star\textquotedblright{}.) Say it once more.
\end{tcolorbox}
\section[\texorpdfstring{\ur{بادل} (bādal)}{بادل (bādal)}]{\ur{بادل} (bādal) — a cloud}

\label{lesson:UR-C65-badal}



\begin{quote}
Before the new one: say the Urdu for a field, then the Urdu for a star.
\end{quote}

\subsection*{You'll want to know: \ur{بادل}}
\textbf{\ur{بادل}} — \emph{bādal} — \textquotedblleft{}a cloud\textquotedblright{}.

\begin{cousinweb}
\textbf{\ur{بادل}} (\emph{bādal}) is a tadbhava, a word inherited from Sanskrit and reshaped by sound change. It comes from \emph{vārdala}, \textquotedblleft{}rainy day, bad weather\textquotedblright{} (CDIAL 11567).
\end{cousinweb}

\begin{grammarlens}[title={What you've built}]
One more word for this chapter.
\end{grammarlens}

\subsection*{Your turn}
\begin{itemize}
\raggedright
  \item \emph{Say it:} \emph{bādal}
  \item \emph{Say it:} \emph{bādal}, once more
  \item \emph{Say it:} say \emph{bādal}
  \item \emph{Recall:} say the Urdu for a field, then the Urdu for a star, then say \emph{bādal} again
\end{itemize}

\begin{tcolorbox}[breakable,colback=teal!4,colframe=teal!35!black,title={Before you move on}]
What does \ur{بادل} mean? (\textquotedblleft{}A cloud\textquotedblright{}.) Say it once more.
\end{tcolorbox}
